\documentclass[12pt,]{article}
\usepackage{lmodern}
\usepackage{amssymb,amsmath}
\usepackage{ifxetex,ifluatex}
\usepackage{fixltx2e} % provides \textsubscript
\ifnum 0\ifxetex 1\fi\ifluatex 1\fi=0 % if pdftex
  \usepackage[T1]{fontenc}
  \usepackage[utf8]{inputenc}
\else % if luatex or xelatex
  \ifxetex
    \usepackage{mathspec}
  \else
    \usepackage{fontspec}
  \fi
  \defaultfontfeatures{Ligatures=TeX,Scale=MatchLowercase}
\fi
% use upquote if available, for straight quotes in verbatim environments
\IfFileExists{upquote.sty}{\usepackage{upquote}}{}
% use microtype if available
\IfFileExists{microtype.sty}{%
\usepackage{microtype}
\UseMicrotypeSet[protrusion]{basicmath} % disable protrusion for tt fonts
}{}
\usepackage[margin=1in]{geometry}
\usepackage[unicode=true]{hyperref}
\hypersetup{
            pdftitle={Lab 1: Math Refresher},
            pdfborder={0 0 0},
            breaklinks=true}
\urlstyle{same}  % don't use monospace font for urls
\usepackage{color}
\usepackage{fancyvrb}
\newcommand{\VerbBar}{|}
\newcommand{\VERB}{\Verb[commandchars=\\\{\}]}
\DefineVerbatimEnvironment{Highlighting}{Verbatim}{commandchars=\\\{\}}
% Add ',fontsize=\small' for more characters per line
\usepackage{framed}
\definecolor{shadecolor}{RGB}{248,248,248}
\newenvironment{Shaded}{\begin{snugshade}}{\end{snugshade}}
\newcommand{\KeywordTok}[1]{\textcolor[rgb]{0.13,0.29,0.53}{\textbf{{#1}}}}
\newcommand{\DataTypeTok}[1]{\textcolor[rgb]{0.13,0.29,0.53}{{#1}}}
\newcommand{\DecValTok}[1]{\textcolor[rgb]{0.00,0.00,0.81}{{#1}}}
\newcommand{\BaseNTok}[1]{\textcolor[rgb]{0.00,0.00,0.81}{{#1}}}
\newcommand{\FloatTok}[1]{\textcolor[rgb]{0.00,0.00,0.81}{{#1}}}
\newcommand{\ConstantTok}[1]{\textcolor[rgb]{0.00,0.00,0.00}{{#1}}}
\newcommand{\CharTok}[1]{\textcolor[rgb]{0.31,0.60,0.02}{{#1}}}
\newcommand{\SpecialCharTok}[1]{\textcolor[rgb]{0.00,0.00,0.00}{{#1}}}
\newcommand{\StringTok}[1]{\textcolor[rgb]{0.31,0.60,0.02}{{#1}}}
\newcommand{\VerbatimStringTok}[1]{\textcolor[rgb]{0.31,0.60,0.02}{{#1}}}
\newcommand{\SpecialStringTok}[1]{\textcolor[rgb]{0.31,0.60,0.02}{{#1}}}
\newcommand{\ImportTok}[1]{{#1}}
\newcommand{\CommentTok}[1]{\textcolor[rgb]{0.56,0.35,0.01}{\textit{{#1}}}}
\newcommand{\DocumentationTok}[1]{\textcolor[rgb]{0.56,0.35,0.01}{\textbf{\textit{{#1}}}}}
\newcommand{\AnnotationTok}[1]{\textcolor[rgb]{0.56,0.35,0.01}{\textbf{\textit{{#1}}}}}
\newcommand{\CommentVarTok}[1]{\textcolor[rgb]{0.56,0.35,0.01}{\textbf{\textit{{#1}}}}}
\newcommand{\OtherTok}[1]{\textcolor[rgb]{0.56,0.35,0.01}{{#1}}}
\newcommand{\FunctionTok}[1]{\textcolor[rgb]{0.00,0.00,0.00}{{#1}}}
\newcommand{\VariableTok}[1]{\textcolor[rgb]{0.00,0.00,0.00}{{#1}}}
\newcommand{\ControlFlowTok}[1]{\textcolor[rgb]{0.13,0.29,0.53}{\textbf{{#1}}}}
\newcommand{\OperatorTok}[1]{\textcolor[rgb]{0.81,0.36,0.00}{\textbf{{#1}}}}
\newcommand{\BuiltInTok}[1]{{#1}}
\newcommand{\ExtensionTok}[1]{{#1}}
\newcommand{\PreprocessorTok}[1]{\textcolor[rgb]{0.56,0.35,0.01}{\textit{{#1}}}}
\newcommand{\AttributeTok}[1]{\textcolor[rgb]{0.77,0.63,0.00}{{#1}}}
\newcommand{\RegionMarkerTok}[1]{{#1}}
\newcommand{\InformationTok}[1]{\textcolor[rgb]{0.56,0.35,0.01}{\textbf{\textit{{#1}}}}}
\newcommand{\WarningTok}[1]{\textcolor[rgb]{0.56,0.35,0.01}{\textbf{\textit{{#1}}}}}
\newcommand{\AlertTok}[1]{\textcolor[rgb]{0.94,0.16,0.16}{{#1}}}
\newcommand{\ErrorTok}[1]{\textcolor[rgb]{0.64,0.00,0.00}{\textbf{{#1}}}}
\newcommand{\NormalTok}[1]{{#1}}
\usepackage{graphicx,grffile}
\makeatletter
\def\maxwidth{\ifdim\Gin@nat@width>\linewidth\linewidth\else\Gin@nat@width\fi}
\def\maxheight{\ifdim\Gin@nat@height>\textheight\textheight\else\Gin@nat@height\fi}
\makeatother
% Scale images if necessary, so that they will not overflow the page
% margins by default, and it is still possible to overwrite the defaults
% using explicit options in \includegraphics[width, height, ...]{}
\setkeys{Gin}{width=\maxwidth,height=\maxheight,keepaspectratio}
\IfFileExists{parskip.sty}{%
\usepackage{parskip}
}{% else
\setlength{\parindent}{0pt}
\setlength{\parskip}{6pt plus 2pt minus 1pt}
}
\setlength{\emergencystretch}{3em}  % prevent overfull lines
\providecommand{\tightlist}{%
  \setlength{\itemsep}{0pt}\setlength{\parskip}{0pt}}
\setcounter{secnumdepth}{0}
% Redefines (sub)paragraphs to behave more like sections
\ifx\paragraph\undefined\else
\let\oldparagraph\paragraph
\renewcommand{\paragraph}[1]{\oldparagraph{#1}\mbox{}}
\fi
\ifx\subparagraph\undefined\else
\let\oldsubparagraph\subparagraph
\renewcommand{\subparagraph}[1]{\oldsubparagraph{#1}\mbox{}}
\fi
\usepackage{enumitem}
\setlist[itemize]{itemsep=4pt, parsep=0pt, topsep=4pt}
\usepackage{etoolbox}
\makeatletter
\providecommand{\subtitle}[1]{% add subtitle to \maketitle
  \apptocmd{\@title}{\par {\large #1 \par}}{}{}
}
\makeatother

\title{Lab 1: Math Refresher}
\providecommand{\subtitle}[1]{}
\subtitle{Instructor notes -- Part 1: prep reading, Part 2: podium/circulating
cheat sheet}
\author{}
\date{\vspace{-2.5em}}

\begin{document}
\maketitle

This single file has two parts, both instructor-only (not a student
handout). \textbf{Part 1} is read beforehand -- full worked examples and
the reasoning behind the flow. \textbf{Part 2} (after the page break,
large print) is the copy to have open while circulating during lab,
since this session isn't one continuous lecture the way M/W periods are
-- it's more likely a short opening refresher followed by spot-checks
with individual groups as they hit each Act.

\section{Where this sits, and why it exists at
all}\label{where-this-sits-and-why-it-exists-at-all}

Lab 1 (Tuesday, Week 1) needs log/exponent rules and a derivative
(product rule + chain rule) in Act 2, and a first taste of integration
in Act 2 Part 3 -- \textbf{before any lecture has formally taught any of
this.} Week 2 Lecture 1 (Monday) is where derivatives, critical points,
and limits get properly developed; that lecture is \emph{next week}.
This doc is what fills the gap: a compressed, rule-focused refresher of
material students are assumed to have seen before (Calc II-ish) but
likely need dusted off, scoped tightly to what Act 2 actually asks for.

\textbf{Deliberately excluded, per the ``hold off on function analysis''
call:} no limits at infinity, no ``throw out the lower-order terms,'' no
critical-point-vs-monotonic vocabulary, no inflection points. All of
that is Week 2 Lecture 1's job, done properly with its own worked
examples (Michaelis-Menten, logistic). Today is purely mechanical
toolkit: the algebra and derivative rules needed to \emph{execute} Act
2, not the conceptual apparatus for \emph{characterizing} a function's
shape.

\section{Rationale for the flow, and an important design
constraint}\label{rationale-for-the-flow-and-an-important-design-constraint}

Order: exponents \(\to\) logarithms (logs are defined in terms of
exponents, so this order is closer to ``necessary'' than ``chosen'')
\(\to\) power rule \(\to\) product rule (both with geometric intuition,
as requested) \(\to\) chain rule \(\to\) quotient rule (derived from
product + chain, a forward-looking bonus since Act 2 doesn't need it)
\(\to\) a first taste of integration, closing with the Fundamental
Theorem of Calculus.

\textbf{Important: every worked numeric example below uses different
numbers than Lab 1's actual questions.} Act 2 gives students a specific
log/exponent warm-up, a specific Ricker function (\(a=0.9\), \(b=1/12\))
to differentiate, and a specific linear-ramp integration problem --
these are graded first-attempt exercises, not lecture content, and Act 2
already hands students the exact rule reminders they need (product rule,
chain rule for \(e^{g(x)}\)) right on the page. Reusing Act 2's actual
numbers here would hand students the answer before they've tried it,
undermining the whole point of the individual-work section. So: the
\emph{rules} and \emph{techniques} below are identical to what Act 2
needs; the specific \emph{numbers and functions} are all different (a
different exponential decay, a different product-times-exponential
function, a different integration ramp) -- same skill, no spoilers. This
is the same design choice made in Week 2 Lecture 1's notes
(Michaelis-Menten instead of Holling IV to teach a technique Lab 2
needed).

\section{Exponent rules}\label{exponent-rules}

\textbf{Core rules} (write these up as a reference block):
\[x^a \cdot x^b = x^{a+b} \qquad \frac{x^a}{x^b}=x^{a-b} \qquad (x^a)^b = x^{ab} \qquad x^{-a}=\frac{1}{x^a}\]

\textbf{Why \(x^0=1\):} don't just assert it -- use the pattern. Take
\(x=2\) and divide by \(x\) each time you drop the exponent by one:
\(2^3, 2^2, 2^1, 2^0 = 8, 4, 2, 1\). Or, more rigorously, from the rule
above: \(x^a/x^a = x^{a-a}=x^0\), and anything divided by itself is
\(1\) (for \(x\neq0\)). Either argument is fine -- the pattern is more
intuitive, the algebra is more rigorous.

\textbf{Worked examples (different numbers than Act 2's):}
\[2^x=32 \Rightarrow x=5 \qquad 5^x=125 \Rightarrow x=3 \qquad 10^x=0.0001 \Rightarrow x=-4\]

For a base that isn't a clean power, e.g. \(e^x=10\), you need
logarithms -- which is exactly the next section.

\section{Logarithm rules}\label{logarithm-rules}

\textbf{What a log actually is, before the rules:} \(\log_b(x)\) is the
exponent you'd raise \(b\) to, to get \(x\). This is worth saying
plainly before any rule -- students who can recite
\(\log(xy)=\log x + \log y\) often can't say what a logarithm \emph{is}.

\textbf{Core rules:}
\[\log(xy)=\log x + \log y \qquad \log(x/y) = \log x - \log y \qquad \log(x^k)=k\log x\]
\[\log_b(b)=1 \qquad \log_b(1)=0 \qquad -\log(x) = \log(1/x)\]

The middle pair are just special cases of the exponent rules above
(since \(b^1=b\) and \(b^0=1\)); the last one follows from
\(\log(x^{-1})=-\log(x)\), itself a special case of
\(\log(x^k)=k\log x\) with \(k=-1\).

\textbf{Order of magnitude -- why \(\log_{10}\) is everywhere in
science:} for a number that's an exact power of ten, \(\log_{10}\) just
counts the zeros: \(\log_{10}(100)=2\), \(\log_{10}(1{,}000{,}000)=6\).
For anything else, \(\log_{10}\) gives you the \emph{order of magnitude}
-- roughly how many digits the number has. E.g.
\(\log_{10}(7{,}000{,}000)\approx6.85\): between \(10^6\) and \(10^7\),
i.e. ``in the millions.'' This is the entire reason log scales exist --
they compress numbers spanning many orders of magnitude (population
sizes, likelihoods, light intensity\ldots{}) onto a scale where equal
\emph{distances} mean equal \emph{multiplicative} factors.

\textbf{Worked examples (different numbers than Act 2's):}
\[\log_{10}(2) + \log_{10}(50) = \log_{10}(100) = 2 \qquad \log_{10}(500)-\log_{10}(5)=\log_{10}(100)=2\]
\[\log_{10}(3^4) = 4\log_{10}(3) \approx 4(0.477) = 1.91\]

\textbf{Negative log = log of the reciprocal, concretely:}
\(-\log_{10}(0.01) = \log_{10}(100) = 2\) -- worth doing side by side
with \(\log_{10}(0.01)=-2\) directly, so students see both routes land
on the same number.

\textbf{Linearizing an exponential decay -- the technique, generically.}
Given \(y=ae^{-bx}\), take \(\ln\) of both sides and apply the product
and power-of-log rules:
\[\ln y = \ln(a\,e^{-bx}) = \ln a + \ln(e^{-bx}) = \ln a - bx\] This is
linear in \(x\): slope \(-b\), intercept \(\ln a\). \textbf{Check with a
numeric example} (again, different from Act 2's \(a=10,\,b=0.5\)):
\(y=5e^{-0.2x}\).

\begin{verbatim}
##    x    y   ln_y
## 1  0 5.00  1.609
## 2  5 1.84  0.609
## 3 10 0.68 -0.391
## 4 15 0.25 -1.391
\end{verbatim}

\(\ln(y)\) drops by exactly \(1.0\) every time \(x\) increases by \(5\)
-- matching \(-b\Delta x = -0.2\times5=-1\) exactly, confirming the
slope interpretation.

\section{Power rule}\label{power-rule}

\textbf{Statement:} \(\dfrac{d}{dx}x^n = n\,x^{n-1}\).

\textbf{Geometric intuition -- \(x^2\) as the area of a growing square.}
Picture a square of side \(x\); area \(x^2\). Grow the side by a small
amount \(dx\). The new area is the old square, plus a thin strip along
the top (\(x \times dx\)), plus a thin strip along the right side
(\(x\times dx\)), plus a tiny corner square (\(dx\times dx = dx^2\))
that's negligible -- it shrinks \emph{faster} than the strips as
\(dx\to0\) (it's second-order small). So
\[\Delta(\text{Area}) \approx x\,dx + x\,dx = 2x\,dx \;\Rightarrow\; \frac{d}{dx}x^2 = 2x\]

\includegraphics{NOTES_Lab1_MathRefresher_files/figure-latex/power-rule-fig-1.pdf}

The same picture generalizes: for \(x^3\) (a cube), growing the side by
\(dx\) adds three faces of area \(x^2\,dx\) each (plus even-smaller
leftover terms), giving \(3x^2\,dx\) -- matching \(n\,x^{n-1}\) with
\(n=3\).

\textbf{Derivation from the product rule, by induction.}
\(x^2 = x\cdot x\); product rule gives \(1\cdot x + x\cdot 1 = 2x\).
\(x^3 = x\cdot x^2\); product rule gives
\(1\cdot x^2 + x\cdot(2x) = 3x^2\) (using the \(x^2\) result just
derived). In general, if \(\frac{d}{dx}x^{n-1}=(n-1)x^{n-2}\) (assume
this holds), then since \(x^n = x\cdot x^{n-1}\):
\[\frac{d}{dx}x^n = 1\cdot x^{n-1} + x\cdot(n-1)x^{n-2} = x^{n-1}+(n-1)x^{n-1} = n\,x^{n-1}\]
Each power's derivative is built from the previous one -- the power rule
is a \emph{consequence} of the product rule, not a separate fact.

\section{Product rule}\label{product-rule}

\textbf{Statement:} \(\dfrac{d}{dx}(uv) = u'v + uv'\).

\textbf{Geometric intuition -- a growing rectangle.} Picture a rectangle
with width \(u(x)\) and height \(v(x)\); area \(uv\). As \(x\) changes
by a small \(dx\), \(u\) changes by \(du=u'dx\) and \(v\) changes by
\(dv=v'dx\) -- generally \emph{different} small amounts, since \(u\) and
\(v\) are different functions. The new area is the old rectangle, plus a
strip along the right (\(v\times du\)), plus a strip along the top
(\(u\times dv\)), plus a tiny corner (\(du\times dv\)) that's negligible
for the same reason as before -- a product of two small things is
smaller still.
\[\Delta(\text{Area}) \approx v\,du + u\,dv \;\Rightarrow\; \frac{d}{dx}(uv) = u'v+uv'\]

\includegraphics{NOTES_Lab1_MathRefresher_files/figure-latex/product-rule-fig-1.pdf}

Same picture as the power rule (unsurprising, since the power rule is a
special case with \(u=v=x\)), just with the two sides allowed to grow at
different rates.

\section{Chain rule}\label{chain-rule}

\textbf{Statement:} \(\dfrac{d}{dx}f(g(x)) = f'(g(x))\cdot g'(x)\), or
in Leibniz notation,
\(\dfrac{dy}{dx} = \dfrac{dy}{du}\cdot\dfrac{du}{dx}\) for \(y=f(u)\),
\(u=g(x)\).

\textbf{Intuition, not geometry this time -- rates multiply through a
composition.} If you bike 3\(\times\) as fast as you walk, and you walk
2\(\times\) as fast as you crawl, then you bike 6\(\times\) as fast as
you crawl -- the rates \emph{multiply}. That's the whole idea:
\(f'(g(x))\) is the outer rate (how fast \(f\) changes per unit of
\(u\)), \(g'(x)\) is the inner rate (how fast \(u\) changes per unit of
\(x\)), and the chain rule multiplies them to get the overall rate.

\textbf{The one form Act 2 actually needs:}
\[\frac{d}{dx}e^{g(x)} = g'(x)\,e^{g(x)}\] (outer function \(f(u)=e^u\)
has \(f'(u)=e^u\); multiply by the inner derivative \(g'(x)\).)

\section{Product + chain together -- a worked example (not
Ricker)}\label{product-chain-together-a-worked-example-not-ricker}

Act 2 asks for the derivative of \(y=a\,x\,e^{-bx}\) using exactly this
combination (product rule across the two factors, chain rule inside the
exponential). Don't work that example here -- it's the actual lab
question. Instead, work a structurally identical warm-up:
\[f(x) = x^2 e^{-x}\]

\begin{Shaded}
\begin{Highlighting}[]
\CommentTok{# symbolic check via numerical differentiation, just to confirm the}
\CommentTok{# hand-derived answer below -- not shown to students, sanity check only}
\NormalTok{f  <-}\StringTok{ }\NormalTok{function(x) x^}\DecValTok{2} \NormalTok{*}\StringTok{ }\KeywordTok{exp}\NormalTok{(-x)}
\NormalTok{fp_formula <-}\StringTok{ }\NormalTok{function(x) x *}\StringTok{ }\NormalTok{(}\DecValTok{2} \NormalTok{-}\StringTok{ }\NormalTok{x) *}\StringTok{ }\KeywordTok{exp}\NormalTok{(-x)}
\NormalTok{x0 <-}\StringTok{ }\FloatTok{1.3}
\NormalTok{(}\KeywordTok{f}\NormalTok{(x0 +}\StringTok{ }\FloatTok{1e-6}\NormalTok{) -}\StringTok{ }\KeywordTok{f}\NormalTok{(x0 -}\StringTok{ }\FloatTok{1e-6}\NormalTok{)) /}\StringTok{ }\FloatTok{2e-6}   \CommentTok{# numerical derivative}
\end{Highlighting}
\end{Shaded}

\begin{verbatim}
## [1] 0.2480039
\end{verbatim}

\begin{Shaded}
\begin{Highlighting}[]
\KeywordTok{fp_formula}\NormalTok{(x0)                          }\CommentTok{# hand-derived formula}
\end{Highlighting}
\end{Shaded}

\begin{verbatim}
## [1] 0.2480039
\end{verbatim}

Product rule with \(u=x^2\) (\(u'=2x\)) and \(v=e^{-x}\)
(\(v'=-e^{-x}\), chain rule with inner derivative \(-1\)):
\[f'(x) = 2x\,e^{-x} + x^2(-e^{-x}) = e^{-x}(2x-x^2) = x(2-x)e^{-x}\]
Critical points at \(x=0\) and \(x=2\) (each factor that can be zero,
set to zero) -- exactly the ``which factor could be zero'' move Act 2's
Step 4 hints at, just demonstrated on a different function first.

\section{Quotient rule -- derived from product + chain (bonus,
forward-looking)}\label{quotient-rule-derived-from-product-chain-bonus-forward-looking}

Not needed for Act 2 (Ricker is a product, not a quotient), but next
week's bestiary is full of quotients (Michaelis-Menten, Holling
II/III/IV) -- worth showing it isn't a fifth rule to memorize, it falls
out of the two already in hand.

Write \(\dfrac{u}{v} = u\cdot v^{-1}\). Apply the product rule:
\[\frac{d}{dx}\left(u\cdot v^{-1}\right) = u'\cdot v^{-1} + u\cdot\frac{d}{dx}\left(v^{-1}\right)\]
The second piece needs the chain rule (outer function \(w^{-1}\), power
rule gives \(\frac{d}{dw}w^{-1}=-w^{-2}\), times inner derivative
\(v'\)): \[\frac{d}{dx}\left(v^{-1}\right) = -v^{-2}\,v'\] Substitute
back and put over a common denominator:
\[\frac{d}{dx}\left(\frac{u}{v}\right) = \frac{u'}{v} - \frac{uv'}{v^2} = \frac{u'v - uv'}{v^2}\]
Every piece of the familiar quotient-rule formula is just product rule +
chain rule + power rule, chained together. Worth saying explicitly: by
the end of today they'll have \emph{all four} differentiation rules
using only two real ideas (product rule, chain rule) plus the algebra
trick \(u/v=uv^{-1}\).

\section{A first taste of
integration}\label{a-first-taste-of-integration}

Two tools in calculus: derivatives (rate of change, from a total) and
\textbf{integrals} (accumulation -- total, from a rate). Act 2 Part 3
already gives students the simplest case fully worked (a constant
capture rate) and asks them to handle a linear ramp themselves -- so
this section reuses their constant-rate example (safe, it's exposition
in the handout, not a graded task) and demonstrates the ramp technique
on \emph{different} numbers before they try their own.

\textbf{Constant rate (Act 2's own worked example, safe to reuse):} rate
\(r=2\) tadpoles/day for \(T=3\) days. Total \(= r\times T = 6\) -- the
area of a \(3\times2\) rectangle.

\textbf{Ramping rate -- worked here with different numbers than Act 2's
``your turn'':} suppose the rate ramps linearly from \(0\) to \(4\)
tadpoles/day over \(T=5\) days.

\includegraphics{NOTES_Lab1_MathRefresher_files/figure-latex/ftc-ramp-fig-1.pdf}

Total = area of the triangle =
\(\frac12\times\text{base}\times\text{height} = \frac12(5)(4)=10\).

\textbf{Fundamental Theorem of Calculus, stated informally:} if \(F\) is
an antiderivative of \(f\) (i.e. \(F'=f\)), then the total accumulated
change from \(a\) to \(b\) is \[\int_a^b f(x)\,dx = F(b)-F(a)\] --
integration and differentiation are inverse operations. Check this
against the triangle above using the power rule \emph{in reverse}: the
ramp is \(f(t) = \frac{4}{5}t\), so an antiderivative is
\(F(t) = \frac{4}{5}\cdot\frac{t^2}{2} = \frac{2}{5}t^2\) (power rule
backwards: raise the power by one, divide by the new power). Then
\[\int_0^5 \frac{4}{5}t\,dt = F(5)-F(0) = \frac{2}{5}(25) - 0 = 10\]
Matches the triangle's area exactly -- the geometric shortcut (area of a
simple shape) and the general machinery (antiderivative, evaluated at
the endpoints) agree, because the geometric shortcut \emph{is} a special
case of the general machinery. This is also a nice full-circle moment:
today opened with the power rule, and it closes by running that same
rule in reverse.

\textbf{Optional aside, only if it lands naturally:} just as a grid
search (Week 1's algorithms lecture) approximates an optimum by checking
many points and gets more exact as you check more of them, a Riemann sum
approximates area by adding up many thin rectangles, and gets more exact
as the rectangles get thinner. Both ideas are ``approximate with many
small pieces, then take a limit'' -- don't force this connection if time
is short, it's a bonus thread, not required.

\section{Take-home message}\label{take-home-message}

\begin{itemize}
\tightlist
\item
  Every rule below is something students have probably seen before;
  today's job is to make it fluent again, not teach it for the first
  time.
\item
  Power rule and product rule both come from the same geometric picture
  (a growing rectangle/square); quotient rule is product + chain rule
  plus one algebra trick (\(u/v=uv^{-1}\)); nothing here is a standalone
  fact to memorize in isolation.
\item
  Every worked example above uses different numbers than Lab 1's actual
  questions -- if a group asks ``is this like the example you showed
  us,'' the honest answer is ``same technique, different numbers, you
  still have to do the algebra.''
\end{itemize}

\section{How to use this during lab}\label{how-to-use-this-during-lab}

This isn't a single continuous lecture -- likely a short (10-15 min)
opening pass through the high points (exponent/log rules, product +
chain rule, one FTC example), then this doc becomes a reference while
circulating. Quick map of what each Act needs, for triage while walking
the room:

\begin{itemize}
\tightlist
\item
  \textbf{Act 2, Part 1} (log/exponent warm-up): exponent rules, log
  rules, linearizing an exponential decay.
\item
  \textbf{Act 2, Part 2} (Ricker derivative + critical point): product
  rule, chain rule (for \(e^{-bx}\)), ``set the derivative to zero and
  solve'' -- the \(x^2e^{-x}\) worked example above is the closest
  parallel to point a stuck group toward.
\item
  \textbf{Act 2, Part 3} (integration taste): the rectangle/triangle
  picture, ``area under a rate curve = total.''
\item
  \textbf{Act 1 and Act 3} don't need any of today's math -- they're
  numeric/visual diagnosis and arithmetic (SSR), respectively. Don't
  spend refresher time there.
\end{itemize}

\newpage

\large

\section{PART 2 -- CHEAT SHEET (circulating / opening
pass)}\label{part-2-cheat-sheet-circulating-opening-pass}

\normalsize
*Copy for the podium or for circulating during lab. Full derivations and
rationale: see Part 1 above.* \large

\subsection{Exponent rules}\label{exponent-rules-1}

\[x^a x^b = x^{a+b} \quad \frac{x^a}{x^b}=x^{a-b} \quad (x^a)^b=x^{ab} \quad x^{-a}=\frac1{x^a} \quad x^0=1\]

\begin{itemize}
\tightlist
\item
  \(x^0=1\): pattern (\(x^3,x^2,x^1,x^0\), dividing by \(x\) each step)
  or \(x^a/x^a=x^{a-a}=x^0=1\)
\item
  \(2^x=32\Rightarrow x=5\); ~ \(10^x=0.0001\Rightarrow x=-4\)
\end{itemize}

\subsection{Logarithm rules}\label{logarithm-rules-1}

\[\log(xy)=\log x+\log y \quad \log(x/y)=\log x-\log y \quad \log(x^k)=k\log x \quad -\log x=\log(1/x)\]

\begin{itemize}
\tightlist
\item
  \(\log_{10}\) = order of magnitude: \(\log_{10}(1{,}000{,}000)=6\)
  (six zeros); non-round numbers land in between
\item
  linearizing decay: \(y=ae^{-bx} \Rightarrow \ln y = \ln a - bx\)
  (line, slope \(-b\), intercept \(\ln a\))
\end{itemize}

\newpage

\subsection{Power rule -- growing
square}\label{power-rule-growing-square}

\[\frac{d}{dx}x^n = n\,x^{n-1}\]

\includegraphics{NOTES_Lab1_MathRefresher_files/figure-latex/power-rule-podium-1.pdf}

\begin{itemize}
\tightlist
\item
  two strips (\(x\,dx\) each) + negligible corner (\(dx^2\))
  \(\to 2x\,dx\)
\item
  from product rule by induction: \(x^n=x\cdot x^{n-1}\to n\,x^{n-1}\)
\end{itemize}

\subsection{Product rule -- growing
rectangle}\label{product-rule-growing-rectangle}

\[\frac{d}{dx}(uv) = u'v + uv'\]

\includegraphics{NOTES_Lab1_MathRefresher_files/figure-latex/product-rule-podium-1.pdf}

\begin{itemize}
\tightlist
\item
  right strip \(v\,du\) + top strip \(u\,dv\) + negligible corner
\end{itemize}

\subsection{Chain rule -- rates
multiply}\label{chain-rule-rates-multiply}

\[\frac{d}{dx}f(g(x)) = f'(g(x))\,g'(x) \qquad\qquad \frac{d}{dx}e^{g(x)} = g'(x)\,e^{g(x)}\]

\begin{itemize}
\tightlist
\item
  bike 3\(\times\) walk, walk 2\(\times\) crawl \(\to\) bike 6\(\times\)
  crawl
\item
  \textbf{warm-up example (not Ricker!):}
  \(f(x)=x^2e^{-x} \to  f'(x)=x(2-x)e^{-x}\), critical pts \(x=0,2\)
\end{itemize}

\subsection{Quotient rule -- product + chain, for
free}\label{quotient-rule-product-chain-for-free}

\[\frac{u}{v}=uv^{-1} \;\to\; \frac{d}{dx}\!\left(\frac{u}{v}\right) = \frac{u'v-uv'}{v^2}\]

\begin{itemize}
\tightlist
\item
  no new idea -- product rule + chain rule (\(v^{-1}\)) + power rule
\item
  previews next week's quotient-shaped bestiary (MM, Holling)
\end{itemize}

\newpage

\subsection{A first taste of
integration}\label{a-first-taste-of-integration-1}

\begin{itemize}
\tightlist
\item
  derivative = rate from total; \textbf{integral} = total from rate
\item
  constant rate \(r=2\), \(T=3\) days \(\to\) total \(=rT=6\)
  (rectangle)
\item
  ramp \(0\to4\) over \(5\) days \(\to\) total \(=\frac12(5)(4)=10\)
  (triangle)
\end{itemize}

\includegraphics{NOTES_Lab1_MathRefresher_files/figure-latex/ftc-podium-1.pdf}

\begin{itemize}
\tightlist
\item
  \textbf{FTC:} \(\int_a^b f(x)dx = F(b)-F(a)\) where \(F'=f\) (inverse
  ops)
\item
  check: \(f(t)=\frac45t\), antiderivative \(F(t)=\frac25t^2\) (power
  rule in reverse) \(\to F(5)-F(0)=10\) -- matches the triangle
\end{itemize}

\subsection{Act-by-Act quick reference (for
circulating)}\label{act-by-act-quick-reference-for-circulating}

\begin{itemize}
\tightlist
\item
  \textbf{Act 2 Pt.1} (log/exp warm-up) -- exponent + log rules
\item
  \textbf{Act 2 Pt.2} (Ricker derivative) -- product rule + chain rule;
  point stuck groups at the \(x^2e^{-x}\) warm-up above, not the answer
\item
  \textbf{Act 2 Pt.3} (integration taste) -- rectangle/triangle, area =
  total
\item
  \textbf{Act 1, Act 3} -- no calculus needed, don't spend time here
\end{itemize}

\end{document}
